\documentclass[11pt]{article}
\usepackage[T1]{fontenc}
\usepackage[utf8]{inputenc}
\usepackage{lmodern}
\usepackage[margin=1in]{geometry}
\usepackage{amsmath,amssymb,amsthm,mathtools}
\usepackage{booktabs,longtable,array}
\usepackage{xcolor}
\usepackage{listings}
\usepackage{microtype}
\usepackage[unicode,hidelinks]{hyperref}
\usepackage{enumitem}
\setlength{\emergencystretch}{3em}
\setlist{itemsep=3pt,topsep=5pt}
\numberwithin{equation}{section}
\newtheorem{theorem}{Theorem}[section]
\newtheorem{lemma}[theorem]{Lemma}
\newtheorem{proposition}[theorem]{Proposition}
\newtheorem{corollary}[theorem]{Corollary}
\theoremstyle{definition}
\newtheorem{definition}[theorem]{Definition}
\newtheorem{remark}[theorem]{Remark}
\newcommand{\R}{\mathbb R}
\newcommand{\Q}{\mathbb Q}
\newcommand{\N}{\mathbb N}
\newcommand{\FR}{\operatorname{FR}}
\newcommand{\calR}{\mathcal R}
\newcommand{\rstar}{r_\star}
\newcommand{\Rzero}{R_0}
\newcommand{\fl}[1]{\left\lfloor#1\right\rfloor}
\newcommand{\ce}[1]{\left\lceil#1\right\rceil}
\newcommand{\normrow}{\mathsf N_D}
\newcommand{\efile}[1]{\texttt{\detokenize{#1}}}
\newcommand{\tagname}[1]{\texttt{\detokenize{#1}}}
\DeclareMathOperator{\conv}{conv}
\lstdefinestyle{code}{
 basicstyle=\ttfamily\scriptsize,
 columns=fullflexible,keepspaces=true,
 showstringspaces=false,breaklines=true,
 breakatwhitespace=false,
 frame=single,framerule=0.3pt,
 rulecolor=\color{black!30},
 aboveskip=9pt,belowskip=9pt,
 numbers=left,numberstyle=\tiny\color{black!55},
 numbersep=7pt,xleftmargin=14pt,
 tabsize=4,upquote=true
}
\lstset{style=code}
\title{A Uniform Lower Bound of \texorpdfstring{$1.802392$}{1.802392}\protect\\
for a General-Bid Factor-Revealing Program}
\author{}
\date{}
\begin{document}
\maketitle
\vspace{-1.5em}
\begin{abstract}
We define a family of factor-revealing maximization problems associated with
continuous, homogeneous, budget-feasible, monotone, superadditive bidding
functions. We prove that, for every admissible bidding function $\beta$,
\[
 \sup_{n\ge1}\FR_n(\beta)>
 \frac{225299}{125000}=1.802392.
\]
Consequently the infimum of these all-size values over admissible bidding
functions is at least $1.802392$. The proof is computer-assisted. It combines
explicit finite feasible configurations, necessary inequalities for bid values,
a finite separation argument, and a rational-arithmetic certificate. The
analytic reductions and exact certificate semantics are given here. The finite
arithmetic calculation relies on the separately supplied certificate
\efile{certificate_1802392.json} and verifier \efile{verify.py}.
No regularity beyond the stated admissibility conditions is assumed, and no
numerical optimizer is needed for verification.
\end{abstract}

\tableofcontents

\section{The optimization problem and the theorem}\label{sec:program}

For $u\in\R$, write $[u]^+=\max\{u,0\}$. All distances and radii below are
nonnegative real numbers; zero distances are permitted.

\begin{definition}[Admissible bidding function]\label{def:bid}
A function $\beta:\R_{\ge0}^2\to\R_{\ge0}$ is \emph{admissible} if it has the
following properties.
\begin{enumerate}[label=(\roman*)]
\item \emph{Nonnegativity:} $\beta(a,d)\ge0$ for $a,d\ge0$.
\item \emph{Budget feasibility:} $\beta(a,d)\le[a-d]^+$ for $a,d\ge0$.
\item \emph{Continuity:} $\beta$ is continuous in both arguments.
\item \emph{Monotonicity:} it is nondecreasing in its first argument and
nonincreasing in its second argument.
\item \emph{Scale equivariance:} for $a>0$ and $d\ge0$,
\[
 \beta(a,d)=a\beta(1,d/a).
\]
\item \emph{Superadditivity along decreasing radii:} whenever
$x_1\ge x_2\ge x_3\ge0$,
\[
 \beta(x_1,x_2)+\beta(x_2,x_3)\le\beta(x_1,x_3).
\]
\end{enumerate}
\end{definition}

Budget feasibility implies $\beta(0,d)=0$ and $\beta(a,a)=0$. Scale equivariance
therefore implies full positive homogeneity,
\begin{equation}\label{eq:homogeneity}
 \beta(ca,cd)=c\beta(a,d)\qquad(c>0,\ a,d\ge0),
\end{equation}
including the cases in which an argument is zero.

\begin{definition}[Factor-revealing program]\label{def:frp}
For an admissible $\beta$ and an integer $n\ge1$, let $\FR_n(\beta)$ be the
supremum of
\begin{equation}\label{eq:objective}
 \sum_{j=1}^{n}\left(\sum_{k=j}^{n}
       \beta(r_{j,k},r_{j,k+1})+r_{j,n+1}\right)-f
\end{equation}
over variables
\[
 f\ge0,\quad d_j\ge0\ (1\le j\le n),\quad
 r_{j,k}\ge0\ (1\le j\le k\le n+1),
\]
subject to
\begin{align}
 \sum_{j=1}^{n}d_j&=1,
 &&\text{normalization},\label{eq:normalization}\\
 r_{j,k+1}&\le r_{j,k},
 &&1\le j\le k\le n,\label{eq:rowmono}\\
 r_{k,k}&\le r_{j,k}+d_j+d_k,
 &&1\le j<k\le n,\label{eq:triangle}\\
 r_{j,\ell}&\le r_{k,\ell}+d_j+d_k,
 &&\substack{1\le j,k<\ell\le n+1\\j\ne k},\label{eq:metric}\\
 \sum_{j=1}^{k-1}\beta(r_{j,k},d_j)
 +\sum_{j=k}^{n}\beta(r_{k,k},d_j)&\le f,
 &&1\le k\le n.\label{eq:facility}
\end{align}
Define the extended-real all-size value
\begin{equation}\label{eq:allsize}
 \calR(\beta):=\sup_{n\ge1}\FR_n(\beta).
\end{equation}
\end{definition}

The terminal terms $r_{j,n+1}$ in \eqref{eq:objective} are part of the
definition. No ordering assumption on the diagonal entries $r_{j,j}$ is added
to the program. Some feasible families used in the proof have ordered
diagonals, which is a restriction on those families only.

\begin{theorem}[Uniform lower bound]\label{thm:main}
Every admissible bidding function satisfies
\begin{equation}\label{eq:main}
 \calR(\beta)>\rstar,
 \qquad \rstar:=\frac{225299}{125000}=1.802392.
\end{equation}
In particular,
\[
 \inf_{\beta\ \mathrm{admissible}}\calR(\beta)\ge1.802392.
\]
\end{theorem}

The strict inequality for each fixed function does not assert a strict
inequality for the infimum. We do not claim that this lower bound is sharp.
The theorem concerns precisely Definition~\ref{def:frp}; it does not assert a
lower bound for any algorithm's actual approximation ratio.

\paragraph{Proof structure and numerical evidence.}
Assume, for contradiction, that $\calR(\beta)\le\rstar$. Many necessary
inequalities will be derived using the weaker bound
\begin{equation}\label{eq:R0}
 \Rzero:=\frac{181}{100}=1.81>\rstar.
\end{equation}
Sections~\ref{sec:bidlemmas}--\ref{sec:rounding} establish the analytic
reductions. Sections~\ref{sec:lift}--\ref{sec:finish} specify the exact finite
proof system and complete the contradiction. The numerical certificate is
supplied separately as \efile{certificate_1802392.json}. Its SHA-256 digest is
\begin{center}\ttfamily\small
4b429ea830a65fcc44fad9974e67f4bebbda\\
005c61644d2cb99768df6df6debc
\end{center}
The line break in the displayed digest is typographic only. The verifier is
supplied separately as \efile{verify.py}; Appendix~\ref{app:reproduce} gives
instructions for running it on the certificate. Its acceptance is a finite
integer/rational calculation, not a numerical-solver infeasibility verdict.

\section{Necessary inequalities for bidding functions}\label{sec:bidlemmas}

Throughout this section suppose $\calR(\beta)\le R$, with $1<R<2$, and set
\begin{equation}\label{eq:gF}
 g(x):=\beta(1,x),\qquad F(t):=\beta(t,1).
\end{equation}
Thus $g$ is nonincreasing, $g(x)=0$ for $x\ge1$, and
$F(t)=t g(1/t)$ for $t>0$.

\subsection{Singletons and the zero-distance normalization}

\begin{lemma}\label{lem:singleton}
For every $t\ge0$,
\begin{equation}\label{eq:singleton}
 F(t)\ge[t-R]^+.
\end{equation}
Moreover,
\begin{equation}\label{eq:zero}
 \beta(a,0)=a\quad(a\ge0),
 \qquad
 [1-Rx]^+\le g(x)\le1-x\quad(0\le x\le1).
\end{equation}
\end{lemma}
\begin{proof}
For $n=1$, set $d_1=1$, $r_{1,1}=r_{1,2}=t$, and $f=F(t)$.
Every constraint holds and the objective is $t-F(t)$. Therefore
$t-F(t)\le R$. Combine this with $F(t)\ge0$ to obtain
\eqref{eq:singleton}. Budget feasibility and scale equivariance imply, for
large $t$,
\[
 1-R/t\le F(t)/t=\beta(1,1/t)\le1-1/t.
\]
Continuity gives $g(0)=1$. Homogeneity then gives $\beta(a,0)=a$; the case
$a=0$ follows from budget feasibility. For $x>0$, applying the singleton
bound at $t=1/x$ gives $g(x)\ge1-Rx$. The remaining assertions follow from
nonnegativity and budget feasibility.
\end{proof}

\begin{lemma}[Positive inverse]\label{lem:inverse}
The function $F$ is continuous and unbounded, and each $v>0$ has a unique
preimage $A(v)>1$ with $F(A(v))=v$. The inverse $A$ is continuous on
$(0,\infty)$.
\end{lemma}
\begin{proof}
Unboundedness follows from \eqref{eq:singleton}. Further,
$F(t)/t=g(1/t)$ is nondecreasing. If $u>t$ and $F(t)>0$, then
$F(u)\ge(u/t)F(t)>F(t)$. Continuity, $F(1)=0$, and unboundedness now give
existence and uniqueness of the positive inverse. Continuity of the inverse
follows from continuity and strict increase on the positive part of the
range.
\end{proof}

\subsection{An outlier construction at every radius}

\begin{lemma}[Outlier inequality]\label{lem:outlier}
For every $s>0$,
\begin{equation}\label{eq:outlier-s}
 s+\beta(s+1,s)-F(s)\le R.
\end{equation}
Equivalently, for $0<x<1$,
\begin{equation}\label{eq:outlier-x}
 g(x)\le R(1-x)-x+xg(1/x-1).
\end{equation}
\end{lemma}
\begin{proof}
Fix $m\ge1$ and put $T=s+1$. Use $n=m+1$ clients, the first with reference
distance $1$ and the remaining $m$ with reference distance $0$. Set
$f=mT+F(s)$. By continuity and monotonicity, there is $a_m\in[s,T]$ such that
\[
 ma_m+F(a_m)=mT+F(s).
\]
Give the first client diagonal $a_m$ and every subsequent entry $s$.
Give each other client every entry $T$. The first facility constraint is the
displayed equality. Every other facility constraint equals
$F(s)+m\beta(T,0)=f$, by Lemma~\ref{lem:singleton}.
The triangle and metric constraints hold since $T-s=1$ is the sum of the
corresponding reference distances. Row monotonicity and normalization also
hold. The objective is
\[
 s+\beta(a_m,s)-F(s).
\]
The defining equation gives
$0\le T-a_m=(F(a_m)-F(s))/m\le(F(T)-F(s))/m$, so $a_m\to T$.
Taking the limit proves \eqref{eq:outlier-s}.
Finally substitute $x=s/(1+s)$, use homogeneity, and multiply by $1-x$ to
obtain \eqref{eq:outlier-x}.
\end{proof}

\subsection{Superadditivity, the golden-ratio bound, and upper envelopes}

\begin{lemma}\label{lem:super}
For $0\le x,y\le1$,
\begin{equation}\label{eq:super}
 g(xy)\ge g(x)+xg(y).
\end{equation}
Consequently, for every integer $k\ge1$,
\begin{equation}\label{eq:power}
 g(x^k)\ge\left(\sum_{p=0}^{k-1}x^p\right)g(x).
\end{equation}
For $0<z<x<1$ one also has
\begin{equation}\label{eq:chain}
 g(z)\ge g(x)+xg(z/x).
\end{equation}
\end{lemma}
\begin{proof}
Apply superadditivity to $1\ge x\ge xy$ and use homogeneity.
If $x=0$, the assertion follows directly from $g(0)=1$.
Equation~\eqref{eq:power} follows by induction, and
\eqref{eq:chain} is \eqref{eq:super} with $y=z/x$.
\end{proof}

Put $\varphi=(1+\sqrt5)/2$. The actual certificate uses $R=\Rzero>\varphi$.

\begin{lemma}\label{lem:golden}
If $R>\varphi$, then
\begin{align}
 g(x)&\le (R-1)(1-x) &&(1/2\le x\le1),\label{eq:halfbnd}\\
 g(x)&\le R-\varphi &&(1/\varphi\le x\le1),\label{eq:goldenbnd}\\
 g(x)&\le R-(1+\varphi)x &&(0\le x<1/\varphi).
 \label{eq:goldenleft}
\end{align}
\end{lemma}
\begin{proof}
For $x\in[1/2,1)$, budget feasibility gives
$g(1/x-1)\le2-1/x$. Substitution in \eqref{eq:outlier-x} yields
\eqref{eq:halfbnd}. At $x=1/\varphi$, the map $x\mapsto1/x-1$ fixes $x$.
Thus \eqref{eq:outlier-x} implies
$(1-x)g(x)\le R(1-x)-x$, or $g(x)\le R-\varphi$.
Monotonicity proves \eqref{eq:goldenbnd} on the entire stated interval.
For $0<x<1/\varphi$, its transformed argument is greater than $1/\varphi$.
If that argument exceeds $1$, its bid is zero; otherwise apply
\eqref{eq:goldenbnd}. In both cases
$g(1/x-1)\le R-\varphi$, since $R>\varphi$. Substitution gives
\eqref{eq:goldenleft}. The endpoints follow directly.
\end{proof}

Here is the exact upper-envelope construction used by the verifier. Set
\begin{equation}\label{eq:philow}
 L=10^{12},\quad k_5=\fl{\sqrt{5L^2}},\quad
 \varphi_-:=\frac{1+k_5/L}{2}
 =\frac{3236067977499}{2000000000000}\le\varphi.
\end{equation}
The integer square-root check $k_5^2\le5L^2<(k_5+1)^2$ certifies the
inequality. The condition $R^2-R-1>0$, together with $R>1$, certifies
$R>\varphi$.

Define $E_0(x)$ to be $1$ at $x=0$, $0$ for $x\ge1$, and, for $0<x<1$, the
minimum of the applicable expressions
\begin{equation}\label{eq:envelope-initial}
\begin{gathered}
 1-x,\qquad (R-1)(1-x)\quad\text{if }x\ge\tfrac12,\\
 R-\varphi_-\quad\text{if }x^2+x\ge1,\qquad
 R-(1+\varphi_-)x\quad\text{if }x^2+x<1.
\end{gathered}
\end{equation}
For $m\ge1$ and $0<x<1$, define recursively
\begin{equation}\label{eq:envelope-rec}
 E_m(x):=\min\{E_0(x),\ R(1-x)-x+xE_{m-1}(1/x-1)\},
\end{equation}
with the same endpoint values. Lemmas~\ref{lem:outlier} and
\ref{lem:golden} imply by induction that $g(x)\le E_m(x)$.
In particular, \eqref{eq:power} gives the further necessary inequality
\begin{equation}\label{eq:envelope-power}
 g(x)\le
 \frac{E_m(x^k)}{\sum_{p=0}^{k-1}x^p}
 \qquad(0<x<1,\ k\ge1).
\end{equation}
All quantities in these recursions are rational for the arguments used in
the certificate. The algorithm uses a finite number of iterations; no
convergence claim is required.

\subsection{Harmonic configurations and certified logarithms}

\begin{lemma}[Harmonic bound]\label{lem:harmonic}
Let $2R-1<T\le3$ and set $u=F(T)$. Then
\begin{equation}\label{eq:harmonic-general}
 R\ge1-p+\frac{(T-1)p}{u}+p\log(u/p)
 \qquad(0<p\le u),
\end{equation}
and
\begin{equation}\label{eq:harmonic-optimized}
 R\ge1+u\exp\bigl((T-1)/u-2\bigr).
\end{equation}
If a rational number $c$ satisfies
\begin{equation}\label{eq:harmonic-test}
 \frac{T-1}{2}<c<T-1,
 \qquad
 \log\frac{c}{R-1}>2-\frac{T-1}{c},
\end{equation}
then $F(T)>c$ and $g(1/T)>c/T$.
\end{lemma}
\begin{proof}
The singleton bound implies $u\ge T-R>(T-1)/2>0$.
Fix $0<p<u$, let $N$ be a positive integer, and set
$M=\fl{Np/u}$. Give all $N$ clients raw reference distance $1$ and choose
raw opening variable $Np$. The last $M$ clients have all entries equal to
$T$. An earlier client with $\ell$ clients remaining has diagonal
$A(Np/\ell)$ and all later entries $1$, where $A$ is the positive inverse
from Lemma~\ref{lem:inverse}. Here $\ell=M+1,\ldots,N$, so its diagonal is
at most $T$ and greater than $1$.

Earlier frozen clients bid $F(1)=0$, and at an early event the remaining
clients' bids sum to $\ell F(A(Np/\ell))=Np$. In the final block, every
client of that block bids $u$, whether active or inactive, so the facility
constraint is $Mu\le Np$. Radii $1$ and $T$ differ by at most $2$, the sum
of two raw reference distances, and all diagonals are at most $T\le3$.
Thus every triangle and metric constraint holds. Divide all variables by
$N$. The resulting objective is
\[
 1-p+(T-1)\frac MN+p\sum_{\ell=M+1}^{N}\frac1\ell.
\]
Letting $N\to\infty$ gives \eqref{eq:harmonic-general}; the endpoint $p=u$
follows by continuity. The harmonic-sum limit can also be obtained by
sandwiching the sum between the integrals of $1/x$ over
$[M,N]$ and $[M+1,N+1]$ for sufficiently large $N$.

The maximizing choice in \eqref{eq:harmonic-general} is
$p=u\exp((T-1)/u-2)$, which belongs to $(0,u)$ because
$u>(T-1)/2$. Substitution gives \eqref{eq:harmonic-optimized}.
The function $w\mapsto1+w\exp((T-1)/w-2)$ is nonincreasing on
$(0,T-1]$, since its derivative has the sign of $1-(T-1)/w$.
Budget feasibility gives $u\le T-1$. If $u\le c$, then
\eqref{eq:harmonic-test} would force the right-hand side of
\eqref{eq:harmonic-optimized} to exceed $R$, a contradiction.
\end{proof}

For completeness, all logarithm comparisons are enclosed as follows.
For $1\le z\le2$, set $y=(z-1)/(z+1)$. Termwise integration of the geometric
series for $(1-y^2)^{-1}$ yields
\begin{equation}\label{eq:logseries}
 \log z=2\sum_{j=0}^{\infty}\frac{y^{2j+1}}{2j+1}.
\end{equation}
The first $n$ terms form a rational lower bound and the remaining tail is
at most
\begin{equation}\label{eq:logtail}
 \frac{2y^{2n+1}}{(2n+1)(1-y^2)}.
\end{equation}
For $z>2$, repeatedly divide by $2$ and add the corresponding multiples of
the interval for $\log2$. For $0<z<1$, negate the interval for $\log(1/z)$.
The verifier uses $n=90$; every comparison in \eqref{eq:harmonic-test} is
made against the lower endpoint of the resulting rational interval.

\section{A finite family of feasible configurations}\label{sec:clients}

\begin{definition}[One-drop client type]\label{def:type}
A type is a triple $z=(d,a,s)$ satisfying
\begin{equation}\label{eq:type}
 d\ge0,\qquad 0\le s\le a\le1+d,\qquad 1-d\le s\le1+d.
\end{equation}
The reference distance is $d$, the diagonal radius is $a$, and the permanent
post-diagonal radius is $s$.
\end{definition}

\begin{lemma}[Geometric feasibility]\label{lem:geometry}
Take a finite multiset of types, order the clients by nondecreasing $a$, and
set, before normalization,
\begin{equation}\label{eq:type-rows}
 d_j=d(z_j),\qquad r_{j,j}=a(z_j),\qquad
 r_{j,k}=s(z_j)\quad(j<k\le n+1).
\end{equation}
All nonnegativity, row monotonicity, triangle, and metric constraints of
Definition~\ref{def:frp} hold. The raw objective contribution of client $j$
is $\beta(a(z_j),s(z_j))+s(z_j)$.
\end{lemma}
\begin{proof}
Nonnegativity and row monotonicity follow from \eqref{eq:type}.
For $j<k$,
\[
 a(z_k)\le1+d(z_k)\le s(z_j)+d(z_j)+d(z_k),
\]
which proves \eqref{eq:triangle}. For any two types,
\[
 s(z_j)-s(z_k)
 \le(1+d(z_j))-(1-d(z_k))=d(z_j)+d(z_k).
\]
Interchanging the two types gives both directed metric inequalities,
including those involving the terminal column. There is only one possible
drop in each row; all later terms have the form $\beta(s,s)=0$.
This proves the objective formula.
\end{proof}

Fix a finite increasing list of activation levels
$\tau_0<\cdots<\tau_{T-1}$, and a finite type list $\mathcal Z$ whose diagonal
radii belong to that list. For $z=(d,\tau_i,s)$, define the true contribution
at a level by
\begin{equation}\label{eq:trueB}
 B^{\rm true}_{t,z}=
 \begin{cases}
  \beta(\tau_t,d),&t\le i,\\
  \beta(s,d),&t>i.
 \end{cases}
\end{equation}
Suppose $C_z\le\beta(\tau_i,s)+s$ and
$B_{t,z}\ge B^{\rm true}_{t,z}$, with $B_{t,z}\ge0$.

\begin{lemma}[Multiplicity bound]\label{lem:multiplicity}
Let $m_z\ge0$ be arbitrary real multiplicities with
$\Delta:=\sum_zd(z)m_z>0$. Then
\begin{equation}\label{eq:multiplicity}
 \frac{\sum_z C_zm_z-\max_t\sum_zB_{t,z}m_z}{\Delta}
 \le\calR(\beta).
\end{equation}
\end{lemma}
\begin{proof}
First suppose all $m_z$ are rational. Multiply by a common denominator and
create that many copies of each type. Order them as in
Lemma~\ref{lem:geometry}. At the beginning of a level $\tau_t$, the
facility contribution is at most $\sum_zB_{t,z}m_z$, up to the common
multiplicity scale. Within a group of clients having the same diagonal,
clients freeze from $\tau_t$ to radii at most $\tau_t$. Monotonicity in the
first argument of $\beta$ implies that the true contribution cannot
increase during this tie group. Thus a constraint at the start of each
level dominates every facility constraint within that level. It does not
matter if some of the listed levels have no clients.

Choose the raw opening variable to be the maximum over the listed levels,
and divide all raw variables by the positive reference-distance sum.
Homogeneity gives a feasible point of Definition~\ref{def:frp}. The lower
objective bound is the left-hand side of \eqref{eq:multiplicity}.
Finally approximate real nonnegative multiplicities by rational ones.
The denominator remains positive, and the numerator is continuous in the
multiplicities because it is a linear function minus the maximum of finitely
many linear functions. Taking limits proves the general case.
\end{proof}

\subsection{One common vector of dual weights}

\begin{lemma}[Finite domination]\label{lem:domination}
Let $B_0,\ldots,B_{T-1}$ and $c$ be vectors in a finite-dimensional real
space. If
\[
 c\cdot m\le\max_t B_t\cdot m\qquad\text{for every }m\ge0,
\]
then there are $\lambda_t\ge0$ with $\sum_t\lambda_t=1$ and
$c\le\sum_t\lambda_t B_t$ coordinatewise.
\end{lemma}
\begin{proof}
Consider the closed convex set
$K=\conv\{B_0,\ldots,B_{T-1}\}-\R_{\ge0}^{|\mathcal Z|}$.
It is closed because it is the sum of a compact set and a closed set.
If $c\notin K$, let $p\in K$ minimize $\|c-p\|$ and put $m=c-p$.
Such a minimum exists by restricting to a sufficiently large closed ball.
For every $z\in K$, differentiating the squared distance along the segment
from $p$ to $z$ at its initial point gives $m\cdot(z-p)\le0$.
Since $K$ is closed under subtracting each nonnegative coordinate direction,
this inequality implies $m\ge0$. But
\[
 m\cdot c=m\cdot p+\|m\|^2>m\cdot p\ge\max_t m\cdot B_t,
\]
contrary to the hypothesis. Thus $c\in K$, which is exactly the stated
coordinatewise domination.
\end{proof}

\begin{proposition}[Necessary client inequalities]\label{prop:dual}
Suppose the type list contains a type with positive reference distance and
$\calR(\beta)\le r$. Then there exists one vector $\lambda$ satisfying
\begin{equation}\label{eq:lambda}
 \lambda_t\ge0,\qquad\sum_{t=0}^{T-1}\lambda_t=1,
\end{equation}
such that, simultaneously for every type $z$,
\begin{equation}\label{eq:clientdual}
 C_z-\sum_t\lambda_tB_{t,z}\le r d(z).
\end{equation}
\end{proposition}
\begin{proof}
Apply Lemma~\ref{lem:multiplicity} and multiply by $\Delta$ to obtain
\[
 \sum_z(C_z-rd(z))m_z\le\max_t\sum_zB_{t,z}m_z
\]
whenever $m\ge0$ and $\Delta>0$. If $\Delta=0$, add a positive
$\varepsilon$ multiple of a positive-distance type and let
$\varepsilon\downarrow0$. This proves the same inequality for every
$m\ge0$. Lemma~\ref{lem:domination}, with $c_z=C_z-rd(z)$, gives the result.
\end{proof}

The type list in this proposition is the \emph{union of all types used
anywhere in the certificate}, not a changing list during verification.
Under the contradictory assumption, choose a single $\lambda$ for this
entire finite union at $r=\rstar$. All earlier deductions at
$\Rzero\ge\rstar$ concern this same $g$ and this same $\lambda$.
This resolves the quantifier issue when additional client rows appear later
in the finite proof.

\section{Finite sampling and one-sided arithmetic}\label{sec:rounding}

\subsection{The exact grids}
Set $q=81/80$. The activation levels are
\begin{equation}\label{eq:timegrid}
 \tau_t=q^t\qquad(0\le t<64).
\end{equation}
The sample grid is the increasing enumeration
$0=x_0<x_1<\cdots<x_J=1$ of
\begin{equation}\label{eq:samplegrid}
 X=\{0\}\ \cup\ \{q^{-k}:0\le k\le153\}
       \ \cup\ \{q^k-1:1\le k\le55\}.
\end{equation}
It has $210$ elements, so $J=209$. The coordinate $g(x_J)=g(1)=0$ is
eliminated. Write
\[
 G_j=g(x_j)\quad(0\le j<J),\qquad G_J:=0.
\]
In particular, $G_0=1$.

For $0\le x<1$, let
\[
 \iota_-(x)=\max\{j:x_j\le x\},\qquad
 \iota_+(x)=\min\{j:x_j\ge x\}.
\]
For $x\ge1$, both indices are defined to be $J$. Monotonicity gives
\begin{equation}\label{eq:sample-sandwich}
 G_{\iota_+(x)}\le g(x)\le G_{\iota_-(x)}.
\end{equation}
Thus the lower and upper sample bids are
\begin{equation}\label{eq:sample-bid}
 \underline\beta(a,d)=
 \begin{cases}aG_{\iota_+(d/a)},&a>d,\\0,&a\le d,\end{cases}
 \quad
 \overline\beta(a,d)=
 \begin{cases}aG_{\iota_-(d/a)},&a>d,\\0,&a\le d.\end{cases}
\end{equation}
They satisfy $\underline\beta\le\beta\le\overline\beta$.
For $d=0<a$ both bounds equal $aG_0=a$.

The client family uses
\begin{align}
 C_z(G)&=s+\underline\beta(a,s),\label{eq:roundedC}\\
 B_{t,z}(G)&=
 \begin{cases}
  \overline\beta(\tau_t,d),&\tau_t\le a,\\
  \overline\beta(s,d),&\tau_t>a.
 \end{cases}\label{eq:roundedB}
\end{align}
These bounds use only neighboring samples; they do not interpolate the
unknown function.

\subsection{Homogeneous row notation and downward rounding}
An affine inequality in the sampled values is written homogeneously as
$a\cdot G\le0$, using $G_0=1$ for its constant term. Any coefficient of
$G_J=0$ is discarded before further operations.

Set
\begin{equation}\label{eq:DS}
 D=S=10^{10}.
\end{equation}
For a finite rational row $a$, define
\begin{equation}\label{eq:normalizer}
 h(a):=\max\{1,\max_j|a_j|\},\qquad
 (\normrow(a))_j:=\fl{D a_j/h(a)}.
\end{equation}
Zero coefficients may be omitted from a stored sparse row.

\begin{lemma}[Safe row rounding]\label{lem:floor}
If $y\ge0$ and $a\cdot y\le0$, then
$\normrow(a)\cdot y\le0$. More generally, replacing any coefficient of a
valid row by a smaller coefficient preserves validity for nonnegative
variables.
\end{lemma}
\begin{proof}
Coordinatewise, $\normrow(a)\le Da/h(a)$. Dotting with $y\ge0$ proves the
claim. This argument applies to negative coefficients as well: floor means
rounding toward $-\infty$, not rounding toward zero.
\end{proof}

The sample bounds in \eqref{eq:sample-sandwich} and the row rounding in
\eqref{eq:normalizer} are distinct operations. Both are one-sided, and neither
uses a numerical tolerance.

\subsection{The initial sampled-bid rows}\label{subsec:rows}
The following real inequalities are valid under $\calR(\beta)\le\Rzero$.
The checker stores their integer normalizations \eqref{eq:normalizer}.
Every subscript $J$ denotes the eliminated zero coordinate.
\begin{align}
 G_i-(1-x_i)G_0&\le0,
 &&\text{\tagname{budget}},\label{eq:tag-budget}\\
 -G_i+(1-\Rzero x_i)G_0&\le0,
 &&\text{\tagname{single}},\label{eq:tag-single}\\
 G_{i+1}-G_i&\le0,
 &&\text{\tagname{monotone}},\label{eq:tag-mon}\\
 G_i+\bigl(x_i-\Rzero(1-x_i)\bigr)G_0
       -x_iG_{\iota_-(1/x_i-1)}&\le0,
 &&\text{\tagname{outlier}}.\label{eq:tag-out}
\end{align}
There are also the following rows, with the indicated side conditions:
\begin{align}
 G_i+(\varphi_- -\Rzero)G_0&\le0,
 &&\text{\tagname{golden}},\label{eq:tag-golden}\\
 G_i+x_iG_j-G_{\iota_-(x_ix_j)}&\le0,
 &&\text{\tagname{super}},\label{eq:tag-super}\\
 \left(\sum_{p=0}^{k-1}x_i^p\right)G_i-G_{\iota_-(x_i^k)}&\le0,
 &&\text{\tagname{power}},\label{eq:tag-power}\\
 -G_i+\frac c{1/x_i}G_0&\le0,
 &&\text{\tagname{harmonic}}.\label{eq:tag-harmonic}
\end{align}
Here \tagname{golden} requires $x_i^2+x_i\ge1$;
\tagname{super} requires $0<x_i,x_j<1$; and \tagname{power} requires
$0<x_i<1$ and $k\ge2$. For \tagname{harmonic}, the side conditions are
exactly Lemma~\ref{lem:harmonic} with $T=1/x_i$ and $R=\Rzero$, including
the certified strict logarithmic comparison. The strict conclusion of that
lemma may safely be weakened to the non-strict row shown here.

The rounded superadditivity rows follow because their positive terms use
exact samples while the negative term uses an upper sample bound. The
outlier row follows in the same way from \eqref{eq:outlier-x}.

The initial tag list in the certificate contains $2201$ rows:
\begin{center}
\begin{tabular}{lr@{\qquad}lr}
\toprule
Tag & Count & Tag & Count\\
\midrule
\tagname{budget}&208&\tagname{single}&208\\
\tagname{monotone}&208&\tagname{outlier}&208\\
\tagname{golden}&55&\tagname{super}&676\\
\tagname{power}&624&\tagname{harmonic}&14\\
\bottomrule
\end{tabular}
\end{center}
The tag records specify the exact indices and, where relevant, the rational
$c$ values. Their validity does not depend on why these particular records
were selected.

\subsection{Additional analytic rows}
The finite proof can append the following further necessary rows.
For $0<u<1$, the tag \tagname{outlier_at} uses
\begin{equation}\label{eq:tag-out-at}
 G_{\iota_+(u)}+\bigl(u-\Rzero(1-u)\bigr)G_0
      -uG_{\iota_-(1/u-1)}\le0.
\end{equation}
For $0<x_k<x_i<1$, the tag \tagname{chain_at} uses
\begin{equation}\label{eq:tag-chain-at}
 G_i+x_iG_{\iota_+(x_k/x_i)}-G_k\le0.
\end{equation}
The positive off-grid terms in these two rows are rounded downward.
The tag \tagname{envelope_power} uses
\begin{equation}\label{eq:tag-env}
 G_i-\frac{E_m(x_i^k)}{\sum_{p=0}^{k-1}x_i^p}G_0\le0
\end{equation}
with the finite recursion \eqref{eq:envelope-rec}. The present certificate
uses $m=12$ for these rows. Finally, the exact rows $-DG_i\le0$ express
nonnegativity, including $-DG_0\le0$.

\section{Lifted variables and a finite proof calculus}\label{sec:lift}

Fix the one common $\lambda$ supplied by Proposition~\ref{prop:dual} for
all $6833$ types at $r=\rstar$. Define
\begin{equation}\label{eq:v}
 v_{t,j}:=\lambda_tG_j,
 \qquad 0\le t<T=64,\quad 0\le j<J=209.
\end{equation}
There are $TJ=13376$ such coordinates. If $e$ is the row with $e_{t,0}=1$
and all other coefficients zero, then
\begin{equation}\label{eq:ebox}
 e\cdot v=1,\quad v\ge0,\quad v_{t,0}=\lambda_t,\quad
 0\le v_{t,j}\le(1-x_j)\lambda_t.
\end{equation}
Furthermore $G_j=\sum_tv_{t,j}$.

The proof only needs to show that every necessary row holds for these
\emph{genuine product variables}. It does not assert that every feasible
point of a lifted linear relaxation admits a factorization
\eqref{eq:v}.

\subsection{Base lifts and client rows}
Every stored integer bid row $a\cdot G\le0$ gives
\begin{equation}\label{eq:base-lift}
 \sum_ja_jv_{t,j}\le0\qquad\text{for each }t.
\end{equation}
For a type $z=(d,\tau_i,s)$, the raw client row is obtained from
\begin{equation}\label{eq:raw-client}
 C_z(G)-\sum_t\lambda_tB_{t,z}(G)-\Rzero d\le0
\end{equation}
by replacing $G_j$ by $\sum_tv_{t,j}$, each $\lambda_tG_j$ by $v_{t,j}$,
and the constant $s-\Rzero d$ by $(s-\Rzero d)e\cdot v$.
The normalized integer row is then stored. Proposition~\ref{prop:dual},
$\rstar\le\Rzero$, and Lemma~\ref{lem:floor} prove its necessity.

\subsection{Deriving and relifting bid inequalities}
Suppose an exact certificate step proves
\begin{equation}\label{eq:derived-g}
 \sum_jc_j\sum_tv_{t,j}\le b.
\end{equation}
For the genuine variables, this is $c\cdot G\le bG_0$.
Thus one may append the bid row
$\normrow(c-be_0)$ and its $T$ separate lifts
\eqref{eq:base-lift}. The same $G$ appears in every block, which is why
relifting a derived inequality is valid.

\subsection{Products with derived dual-weight inequalities}
Suppose an earlier certified step proves
\begin{equation}\label{eq:lambda-cut}
 \sum_th_t\lambda_t\le b.
\end{equation}
For any already valid stored bid row $a\cdot G\le0$,
\begin{equation}\label{eq:product-exact}
 \sum_t(b-h_t)\sum_ja_jv_{t,j}
 =\left(b-\sum_th_t\lambda_t\right)(a\cdot G)\le0.
\end{equation}
Therefore the integer row
\begin{equation}\label{eq:product-floor}
 P_{t,j}=\fl{(b-h_t)a_j}
\end{equation}
is valid. This last rounding is directly applied to the already-integer
$a_j$; it does not apply \eqref{eq:normalizer} a second time.

Individual $b-h_t$ need not be nonnegative. The proof is the scalar product
identity \eqref{eq:product-exact}, followed by coefficientwise downward
rounding. It would not be valid to justify the row by treating each block
multiplier as separately nonnegative.

\subsection{Strengthening client rows to the final threshold}
Let $p_z$ be the raw rational coefficient vector in
\eqref{eq:raw-client}, and let $h_z=h(p_z)$ be its normalization scale.
The raw row at a smaller threshold $R'$, with
$\rstar\le R'\le\Rzero$, is
\[
 p_z+(\Rzero-R')d(z)e.
\]
If the original integer row is $A_z=\normrow(p_z)$, the verifier appends
\begin{equation}\label{eq:refine}
 A'_z=A_z+\nu_ze,
 \qquad
 \nu_z=\fl{D(\Rzero-R')d(z)/h_z}.
\end{equation}
This preserves the \emph{original} scale $h_z$.
For constant coordinates, $\lfloor u\rfloor+\lfloor w\rfloor\le u+w$;
for all other coordinates the same base rounding is retained. Hence
\[
 A'_z\le \frac D{h_z}
       \bigl(p_z+(\Rzero-R')d(z)e\bigr)
\]
coordinatewise. Since the raw row is necessary at $R'$ for the common
$\lambda$ chosen at $\rstar$, \eqref{eq:refine} is valid.
The final certificate applies this rule to $400$ existing client rows, with
$R'=\rstar$. No new geometric configuration is asserted at this step.

\section{Exact checking of derived inequalities}\label{sec:arithmetic}

All global rows have integer coefficients and the orientation
$A_i\cdot v\le0$. Row indices refer only to rows already available at the
step being checked.

For $j>0$, put
\begin{equation}\label{eq:Uj}
 U_j:=\ce{D(1-x_j)}.
\end{equation}
Then \eqref{eq:ebox} implies
$v_{t,j}\le(U_j/D)\lambda_t$.

\begin{lemma}[Budget-aware residual certificate]\label{lem:residual}
Let $c$ be a rational objective row for which every $SDc_{t,j}$ is an
integer. Choose nonnegative integers $m_i$, a signed integer $\mu$, and
define the integer residual
\begin{equation}\label{eq:Q}
 Q_{t,j}:=\sum_im_iA_{i,t,j}+\mu D e_{t,j}-SDc_{t,j}.
\end{equation}
Set
\begin{equation}\label{eq:error}
 E:=\max_{0\le t<T}\left\{
 -DQ_{t,0}+\sum_{j=1}^{J-1}\max\{-Q_{t,j},0\}\,U_j
 \right\}.
\end{equation}
For every genuine $v$ satisfying the available rows,
\begin{equation}\label{eq:witnessbound}
 c\cdot v\le H:=\frac{\mu D^2+E}{SD^2}.
\end{equation}
\end{lemma}
\begin{proof}
Nonnegativity of the multipliers gives
\[
 Q\cdot v
 =\sum_im_iA_i\cdot v+\mu D-SD(c\cdot v)
 \le\mu D-SD(c\cdot v).
\]
For each block, the bounds \eqref{eq:ebox} and \eqref{eq:Uj} give
\[
 -\sum_jQ_{t,j}v_{t,j}
 \le\left(-Q_{t,0}+\sum_{j>0}
       \max\{-Q_{t,j},0\}\frac{U_j}{D}\right)\lambda_t.
\]
Sum over $t$ and use $\sum_t\lambda_t=1$ to obtain
$-Q\cdot v\le E/D$. Substitution and division by $SD$ give
\eqref{eq:witnessbound}.
\end{proof}

A witness records its objective $c$, its nonnegative integer multipliers,
$\mu$, and a rational \emph{reported upper bound} $b$. The checker verifies
exactly that the recorded objective is the expected one and that $H\le b$.
It then uses $c\cdot v\le b$. In particular, numerical residuals are never
discarded; \eqref{eq:error} charges every one in the conservative direction.

For a bid-cut step, $c_{t,j}$ is independent of $t$ and equals the requested
bid direction. For a dual-weight-cut step, $c_{t,0}=h_t$ and all other
coefficients are zero. For the terminal contradiction, $c=0$.

\section{The finite certificate and completion of the proof}\label{sec:finish}

\subsection{Data and deterministic reconstruction}
The JSON object contains exact rational strings for grid values, distances,
radii, directions, and reported bounds. Integer fields specify indices,
scales, and witness multipliers. The essential fields are:
\begin{center}
\begin{tabular}{@{}ll@{}}
\toprule
Field & Meaning\\
\midrule
\tagname{kind} & The value \tagname{geom}\\
\tagname{R}, \tagname{target_R} & $\Rzero=181/100$, $\rstar=225299/125000$\\
\tagname{D}, \tagname{S} & The two exact scales $10^{10}$\\
\tagname{xs}, \tagname{times} & The grids \eqref{eq:samplegrid}, \eqref{eq:timegrid}\\
\tagname{gtags}, \tagname{base_grows} & The initial $2201$ bid-row tags\\
\tagname{types} & $1067$ initial records $(i,d,s)$ with $a=\tau_i$\\
\tagname{proof} & The ordered list of $156$ proof events\\
\bottomrule
\end{tabular}
\end{center}
A later \tagname{clients} event adds $5766$ types. These and the initial
ones are $6833$ distinct geometries, each checked against
\eqref{eq:type}. The list includes one zero-reference-distance type and at
least one positive-reference-distance type.

All indices are zero-based. Initially the global rows are ordered by first
listing the $64$ lifts of bid row $0$, then the $64$ lifts of bid row $1$,
and so on, followed by the $1067$ initial client rows. Thus the starting
number of global rows is
\[
 2201\cdot64+1067=141931.
\]
All later rows are appended in the order specified in the event list.
Appendix~\ref{app:format} specifies the row numbering rules; the separately
supplied \efile{verify.py} implements them.

The proof-event counts are
\begin{center}
\begin{tabular}{@{}lrp{7.3cm}@{}}
\toprule
Event & Count & Effect\\
\midrule
\tagname{axiom}&1&Store the $209$ nonnegativity bid rows.\\
\tagname{axiom_g}&15&Append $1419$ analytic bid rows and their lifts.\\
\tagname{envelope_axioms}&1&Append $208$ envelope-power bid rows and their lifts.\\
\tagname{g}&7&Verify, store, and lift a derived bid inequality.\\
\tagname{lambda}&14&Verify a dual-weight inequality and append designated products.\\
\tagname{clients}&1&Append the $5766$ additional geometric client rows.\\
\tagname{product}&115&Append further products using earlier dual-weight cuts.\\
\tagname{refined_clients}&1&Tighten $400$ client rows to $\rstar$.\\
\tagname{contradiction}&1&Verify a strictly negative upper bound on the zero objective.\\
\bottomrule
\end{tabular}
\end{center}
The completed system has $4044$ stored bid rows and $259756$ global rows.
There are $22$ arithmetic witnesses: $7$ bid cuts, $14$ dual-weight cuts,
and the terminal contradiction.

\subsection{Soundness of the proof interpreter}

\begin{proposition}[Certificate soundness]\label{prop:soundness}
Under the assumption $\calR(\beta)\le\rstar$, every global row and every
stored bid inequality reconstructed by the verifier is valid for the genuine
variables \eqref{eq:v}. Every accepted witness gives the asserted upper
bound. Therefore acceptance of the terminal negative-bound witness is a
contradiction.
\end{proposition}
\begin{proof}
Choose the common $\lambda$ for the union of all client types as specified
after Proposition~\ref{prop:dual}. The initial bid rows are valid by
Section~\ref{sec:bidlemmas}, the sample sandwich, and
Lemma~\ref{lem:floor}. Their lifts are valid because $\lambda_t\ge0$.
The initial client rows are valid by Proposition~\ref{prop:dual} and
one-sided rounding.

Induct on the proof-event index. An analytic-axiom event is valid by
\eqref{eq:tag-out-at}, \eqref{eq:tag-chain-at},
\eqref{eq:tag-super}, or \eqref{eq:tag-env}, as applicable.
A client event is valid by its explicitly checked geometry and
Proposition~\ref{prop:dual}. A derived bid inequality follows from
Lemma~\ref{lem:residual}, then may be relifted as in
\eqref{eq:derived-g}. A derived dual-weight inequality also follows from
Lemma~\ref{lem:residual}; its designated products are justified by
\eqref{eq:product-exact}--\eqref{eq:product-floor}. A subsequent product
event is identical, except that its dual-weight inequality was proved at an
earlier event. The interpreter checks these earlier-index requirements.
A client refinement is valid by \eqref{eq:refine}.
Thus every row available to the next witness is necessary for the genuine
variables. Finally a witness with $c=0$ and reported bound $b<0$ gives
$0\le b<0$, which is impossible.
\end{proof}

\subsection{The checked arithmetic result}
The final witness uses $12484$ multiplier records and the signed integer
\[
 \mu=-24859.
\]
Exact evaluation of \eqref{eq:Q}--\eqref{eq:witnessbound} verifies a bound
no larger than the recorded value
\begin{equation}\label{eq:negative}
 b_{\rm final}=-\frac{1241}{500000000}.
\end{equation}
Its objective list is empty, meaning $c=0$. Thus its accepted assertion is
\[
 0\le-\frac{1241}{500000000}<0.
\]
The separately supplied verifier \efile{verify.py} reconstructs the raw axioms,
grid, client rows, normalizations, intermediate inequalities, and final
residual correction.
It checks the complete certificate in one run and does not use an
optimization solver.

\begin{proof}[Proof of Theorem~\ref{thm:main}]
If $\calR(\beta)=+\infty$, the conclusion is immediate. Otherwise suppose
$\calR(\beta)\le\rstar$. All reductions and necessary inequalities proved
above apply. The separately supplied \efile{certificate_1802392.json} is
accepted by \efile{verify.py}, which implements the exact checks of
Sections~\ref{sec:lift}--\ref{sec:finish} and Appendix~\ref{app:format},
with final bound \eqref{eq:negative}. The reproduction command is given in
Appendix~\ref{app:reproduce}. Proposition~\ref{prop:soundness} therefore gives a
contradiction. Hence $\calR(\beta)>\rstar$ for every admissible $\beta$.
Taking the infimum proves the stated non-strict minimax lower bound.
\end{proof}

\paragraph{What the calculation establishes.}
The finite data are essential to the computer-assisted portion of the proof.
The analytic part proves why acceptance of those data implies the theorem;
the verifier supplies the reproducible finite arithmetic calculation.
This is not a proof-assistant formalization. No claim is made that a
particular bid is extremal, that the sample system exactly describes all
admissible bids, or that $\rstar$ is the best possible lower bound.

\appendix
\section{Certificate format and indexing details}\label{app:format}

This appendix specifies the event semantics sufficiently to reimplement the
checker. The executable implementation is supplied separately as
\efile{verify.py}.
A sparse row is a map from coordinate indices to rational or integer
coefficients; coefficients at the same index are added before normalization.
A lifted coordinate $(t,j)$ is flattened as $tJ+j$.

\subsection{Initial tag records}
The initial records have one of the following forms:
\begin{center}
\begin{tabular}{@{}ll@{}}
\toprule
Record & Defining inequality\\
\midrule
\tagname{[budget,i]} & \eqref{eq:tag-budget}\\
\tagname{[single,i]} & \eqref{eq:tag-single}\\
\tagname{[monotone,i]} & \eqref{eq:tag-mon}\\
\tagname{[outlier,i]} & \eqref{eq:tag-out}\\
\tagname{[golden,i]} & \eqref{eq:tag-golden}\\
\tagname{[super,i,j]} & \eqref{eq:tag-super}\\
\tagname{[power,i,k]} & \eqref{eq:tag-power}\\
\tagname{[harmonic,i,c]} & \eqref{eq:tag-harmonic}, with \eqref{eq:harmonic-test}\\
\bottomrule
\end{tabular}
\end{center}
Later analytic tags additionally have the forms
\begin{center}
\tagname{[outlier_at,u]},\qquad \tagname{[chain_at,i,k]},\\[3pt]
\tagname{[envelope_power,i,k,m]}.
\end{center}
Their rational arguments and indices are part of the certificate, not
floating-point inputs. The verifier checks their side conditions.

\subsection{Witness records}
A witness has fields \tagname{mode}, \tagname{objective},
\tagname{multipliers}, \tagname{mu}, and \tagname{upper}.
Here \tagname{mode} is \tagname{budget}; \tagname{objective} is a list of
pairs $(\text{coordinate},\text{rational coefficient})$;
\tagname{multipliers} is a list of pairs
$(\text{already available row index},\text{nonnegative integer})$;
\tagname{mu} is a signed integer; and \tagname{upper} is the reported rational
bound. Equation~\eqref{eq:Q} must be integral at every coordinate.
The witness objective must exactly match the objective prescribed by its
containing event.

\subsection{Events}
The \tagname{axiom} event names \tagname{nonnegative_g} and gives coordinate
indices. It appends $-De_j$ to the stored bid-row list, but appends no global
rows. These bid rows can later be used in products. Nonnegativity of $v$ is
already part of the mathematical invariant \eqref{eq:ebox}.

An \tagname{axiom_g} or \tagname{envelope_axioms} event contains a list of
analytic tags. For each tag in order, its normalized row is appended to the
bid-row list, then its $T$ lifts are appended to the global row list in
increasing time order.

A \tagname{clients} event contains a list of $(i,d,s)$ records. For each in
order, check \eqref{eq:type} with $a=\tau_i$, build
\eqref{eq:raw-client}, and append its normalized row.

A \tagname{g} event gives a sparse direction $c$, a bound $b$, and a witness.
Its lifted objective repeats $c$ identically in every time block. Check the
witness, check that its reported bound equals $b$, append
$\normrow(c-be_0)$ to the bid-row list, and append all its time-block lifts.

A \tagname{lambda} event gives a sparse direction $h$, a bound $b$, a
witness, and a list of stored bid-row indices. Its objective is supported on
coordinates $(t,0)$ with coefficients $h_t$. After verifying the witness and
bound, append the product row \eqref{eq:product-floor} for each designated
bid-row index, in list order. The bound $(h,b)$ is also available for later
product events. There is no extra direct global row for the weight
inequality unless a requested product supplies it.

A \tagname{product} event identifies an earlier \tagname{lambda} event by
its zero-based proof-event index and lists bid-row indices that already
exist. Append their product rows using that earlier event's $(h,b)$.

A \tagname{refined_clients} event supplies $R'$ and a list of original
client-row indices. Check $\rstar\le R'\le\Rzero$, reconstruct each source
row's original normalization scale, and append \eqref{eq:refine}. The
present certificate uses only ordinary client rows as sources.

Finally, \tagname{contradiction} contains a witness for the empty objective.
It must have negative reported bound. The checker rejects a file without a
terminal contradiction and rejects any event after that contradiction.

\section{Reproduction}\label{app:reproduce}

\subsection{Certificate verification}
The computer-assisted calculation requires the separately supplied files
\efile{certificate_1802392.json} and \efile{verify.py}. The verifier implements
the proof calculus and exact arithmetic checks described in
Sections~\ref{sec:lift}--\ref{sec:finish} and Appendix~\ref{app:format}.
It is specialized to the stated grids, constants, and certificate and uses
only Python's standard library. Integer and rational operations decide every
assertion; a floating-point clock is used only to display elapsed time.

Place both files in one directory and use Python 3.10 or later with assertions
enabled. From that directory, run
\begin{lstlisting}[language=bash,numbers=none]
python verify.py certificate_1802392.json
\end{lstlisting}
Do not use Python's optimization flag or enable optimization through
\texttt{PYTHONOPTIMIZE}; the executable rejects execution with assertions
disabled. The expected final output includes
\begin{lstlisting}[numbers=none]
COMPLETE EXACT CERTIFICATE VERIFIED.
Assumed upper bound: 225299/125000
Impossible conclusion: 0 <= -1241/500000000
Therefore every admissible beta has sup_n FR_n(beta) > 225299/125000
\end{lstlisting}
The verification log is supplementary evidence, not an input to the proof.
The numerical certificate is checked from its contents on every run.

\subsection{Compilation}
The certificate and verifier are required for the arithmetic verification,
while this TeX file compiles without them, bibliography files, or image files.
A standard LaTeX installation with the packages named in its preamble suffices.
Run
\begin{lstlisting}[language=bash,numbers=none]
pdflatex -interaction=nonstopmode -halt-on-error frp_lower_bound_1802392.tex
pdflatex -interaction=nonstopmode -halt-on-error frp_lower_bound_1802392.tex
\end{lstlisting}
The later passes resolve the table of contents and internal cross-references.
Repeat the second command if LaTeX requests another pass.

\end{document}
